%=============================================================================
\section{Experimental Setup}
\label{sec:setup}
%=============================================================================

Every number reported in this paper comes from the same harness, described in
Section~\ref{sec:system}, run under the protocol fixed here. The protocol is
stated in full before any result is shown because several of the choices---the
budget, the paired sampling, the separation of assumed from true
reliability---determine what the results can and cannot be read to mean.

\subsection{Catalogues}
\label{sec:catalogues}

We use three real catalogues from three domains and one synthetic family.
Table~\ref{tab:catalogues} summarises them.

\emph{Music.} A catalogue of $79{,}803$ tracks described by $21$ attributes,
derived from a public audio-features corpus and used as the item table of a
deployed Akinator-style music identification system. Eleven of the attributes
are discretisations of continuous audio descriptors (danceability, energy,
valence, acousticness, instrumentalness, liveness, speechiness, loudness,
tempo, popularity, duration); the remainder are natively categorical metadata
(genre, language, content rating, artist type, release type, track version) or
music-theoretic labels (key, mode, time signature). This mixture is the
structure the mechanism of Section~\ref{sec:discretisation-theory} requires,
and it is why this catalogue is the primary one.

\emph{Mushroom.} UCI Agaricus-Lepiota~\cite{uci}, $8{,}124$ specimens by $21$ usable
categorical attributes, cast as field identification by a forager. Every
attribute is natively categorical, so the discretisation mechanism cannot
operate. The catalogue is nonetheless a genuine test rather than a formality,
because its two most informative attributes---odour, and colour names drawn
from nine to twelve options---are the two classic unreliable human perceptual
judgements. If high information gain concentrated on hard-to-answer attributes
for reasons other than binning, this is where it would show.

\emph{Dermatology.} UCI Dermatology~\cite{uci,guvenir1998}, $366$ cases by $34$ attributes, cast as
differential narrowing by a triage questioner with the patient as oracle. This
is the one catalogue whose reliability annotation we did not invent: the
dataset documentation itself partitions attributes into clinical and
histopathological, and a histopathological attribute is by construction one the
patient cannot answer at all, since it requires a biopsy and a microscope. Our
annotation refines that documented split by asking who can answer from what
evidence. The single continuous column, age, is banded into quartiles.

\emph{Synthetic.} A family in which the rank correlation between an attribute's
apparent informativeness and its reliability is a control knob, used in
Section~\ref{sec:results-synthetic} to turn the randomised and adversarial
controls---which are two isolated points on that axis---into a continuum.

\begin{table}[t]
\caption{Catalogues. $|\Qset|$ is the number of compiled one-versus-rest
questions after degenerate questions are dropped. The entropy floor is
$\log_2$ of the number of distinct attribute vectors, i.e.\ the information
that must be acquired even with perfect answers; $T_{\min}$ is the
capacity-corrected budget bound~\eqref{eq:budget} under each catalogue's
reliability annotation.}
\label{tab:catalogues}
\centering
\footnotesize
\begin{tabular}{@{}lrrrrr@{}}
\toprule
\textbf{Catalogue} & \textbf{$N$} & \textbf{$M$} & \textbf{$|\Qset|$} &
\textbf{floor} & \textbf{$T_{\min}$} \\
\midrule
Music        & 79{,}803 & 21 & 193 & 16.23 & 20.15 \\
Mushroom     & 8{,}124  & 21 & 116 & 12.99 & 16.12 \\
Dermatology  & 366      & 34 & 134 & 8.52  & 10.58 \\
\bottomrule
\end{tabular}
\end{table}

The music catalogue is not perfectly identifiable, and this bounds every
accuracy we report. Its $79{,}803$ tracks occupy $77{,}069$ distinct attribute
vectors, so $93.7\%$ of items are uniquely identifiable and the largest
equivalence class holds $11$ indistinguishable tracks. No method, with any
budget and a perfect user, can exceed $96.6\%$ top-1 accuracy on it; the clean
condition confirms that the best methods reach exactly that ceiling.

\subsection{Attribute Inventory}
\label{sec:attributes}

Table~\ref{tab:attributes} lists the music attributes with their annotated
tier, their maximum one-versus-rest information gain in the catalogue as
shipped, and the same quantity after the eleven continuous sources are
re-discretised into four equal-population quantile bins. The second column is
the operative one for Section~\ref{sec:results-bias}; the third exists because
it is the clean instantiation of Proposition~\ref{prop:discretisation}, and it
is where every re-binned attribute reports exactly $\Hb(1/4)=0.811$ bits,
regardless of what the attribute means.

Two features of the table should be noted in advance. First, the shipped
catalogue's continuous attributes are \emph{not} equal-population bins: they
were cut at fixed semantic thresholds by the original pipeline, so their split
masses vary and their maximum information gain ranges from $0.231$ to $1.000$
rather than sitting at $\Hb(1/5)$. The bias in the shipped catalogue is
therefore not an artefact of our own preprocessing---we did not do the
binning---and the four-bin column shows what happens when the binning is made
textbook-clean, which is that the bias strengthens
(Section~\ref{sec:results-bias}). Second, genre carries $113$ values and a
maximum one-versus-rest information gain of $0.097$ bits: it is the most
discriminative attribute in the catalogue by any reasonable measure and the
least attractive to a criterion that maximises binary entropy. That single row
is the whole problem in miniature.

\begin{table}[t]
\caption{Music attributes. $|V|$ is the number of distinct values;
$\max\Hb$ is the largest one-versus-rest binary entropy at a uniform belief;
$\max\Hb^{(4)}$ is the same after re-discretising the underlying continuous
column into four equal-population bins (``---'' where the attribute has no
continuous source). Release type is degenerate in this catalogue (a single
value) and compiles to no questions.}
\label{tab:attributes}
\centering
\footnotesize
\begin{tabular}{@{}llrrr@{}}
\toprule
\textbf{Attribute} & \textbf{Tier} & \textbf{$|V|$} &
\textbf{$\max\Hb$} & \textbf{$\max\Hb^{(4)}$} \\
\midrule
genre            & objective & 113 & 0.097 & --- \\
language         & objective & 2   & 0.446 & --- \\
content rating   & objective & 2   & 0.424 & --- \\
artist type      & objective & 4   & 0.834 & --- \\
track version    & objective & 5   & 0.424 & --- \\
release type     & objective & 1   & ---   & --- \\
\midrule
popularity       & semi & 5 & 0.958 & 0.811 \\
duration         & semi & 5 & 0.939 & 0.811 \\
acousticness     & semi & 5 & 0.999 & 0.811 \\
instrumentalness & semi & 5 & 0.819 & 0.811 \\
liveness         & semi & 5 & 0.933 & 0.811 \\
speechiness      & semi & 3 & 0.231 & 0.811 \\
loudness         & semi & 5 & 1.000 & 0.811 \\
\midrule
mood             & subjective & 5 & 1.000 & --- \\
tempo            & subjective & 5 & 0.902 & 0.811 \\
danceability     & subjective & 5 & 0.957 & 0.811 \\
energy           & subjective & 5 & 0.924 & 0.811 \\
valence          & subjective & 5 & 0.804 & 0.811 \\
key              & subjective & 4 & 0.896 & --- \\
mode             & subjective & 2 & 0.947 & --- \\
time signature   & subjective & 3 & 0.496 & --- \\
\bottomrule
\end{tabular}
\end{table}

\subsection{Baselines}
\label{sec:baselines}

A comparison of ``information gain'' against ``a learned policy'' confounds two
independent choices: how the next question is selected, and how the answer is
absorbed. Proposition~\ref{prop:limit} shows the second choice is a limit of one
model rather than a different model, so the two can and should be crossed. Our
nine methods do that, and are listed in Table~\ref{tab:baselines}.

\begin{table}[t]
\caption{Methods. ``Selector'' is the rule that chooses the next question;
``update'' is how the answer is absorbed; $\hat\varepsilon$ is what the method
assumes about answer reliability. Hard elimination is
$\hat\varepsilon \equiv 0$ (Proposition~\ref{prop:limit}), so the selector and
update columns are not independent for that row alone.}
\label{tab:baselines}
\centering
\scriptsize
\begin{tabular}{@{}llll@{}}
\toprule
\textbf{Method} & \textbf{Selector} & \textbf{Update} & \textbf{$\hat\varepsilon$} \\
\midrule
Random+soft         & uniform random & soft & uniform \\
IG+hard             & $\Hb(p_q)$     & hard & $0$ \\
IG+soft-uniform     & $\Hb(p_q)$     & soft & uniform \\
IG+soft-objective   & $\Hb(p_q)$, objective only & soft & uniform \\
Learned-noiseblind  & ridge value regression & soft & uniform \\
Learned-noiseaware  & ridge value regression & soft & uniform \\
RAIG-empirical      & $\RAIG(q)$     & soft & measured \eqref{eq:discordance} \\
RAIG-tiered         & $\RAIG(q)$     & soft & 3-tier annotation \\
RAIG-oracle         & $\RAIG(q)$     & soft & true $\varepsilon$ \\
\bottomrule
\end{tabular}
\end{table}

Two further specifications of $\hat\varepsilon$---RAIG-EM, which learns the
crossover vector from interaction logs, and RAIG-psycho, which computes it from
the catalogue's own cut points---are introduced where they are tested
(Sections~\ref{sec:results-em} and~\ref{sec:results-psycho}) rather than here,
because each is evaluated against this nine-method set rather than alongside it.

Three of the nine deserve comment.

\emph{IG+soft-uniform is the honest incumbent.} It is classical information
gain paired with a correct probabilistic belief update at the catalogue-average
crossover. By Proposition~\ref{prop:uniform} its selector is exactly classical
information gain, so the comparison of RAIG-tiered against IG+soft-uniform
isolates the acquisition function with the update held fixed. This is the
pre-registered primary comparison, and the only one we treat as confirmatory.

\emph{IG+soft-objective is the restricted-attribute baseline}, the practitioner
heuristic of simply not asking unreliable questions. It is the baseline that
most directly threatens the method---if discarding subjective attributes
worked, weighting them would be unnecessary---and we report it in every
condition. Its weakness is structural rather than empirical: the six objective
attributes partition the music catalogue into only $1{,}926$ equivalence
classes, the largest holding $932$ tracks, so the strategy has a top-1 accuracy
ceiling of $0.65\%$ however many questions it is given. We state this ceiling
in advance so that its poor showing is not mistaken for a tuning failure.

\emph{The learned policies} are myopic value-regression selectors in the
EAR/CPR mould~\cite{lei2020ear,lei2020cpr}: roll out a behaviour policy, log
state--question features against the realised one-turn reduction in belief
entropy, fit a ridge regressor, and select by the learned score. Training uses
$600$ simulated sessions of $T=20$ turns, giving $12{,}000$ samples over $51$
features, with ridge penalty $\lambda=1$. The two variants differ in exactly
one respect: the noise used to generate the training rollouts. Noiseblind
trains on noiseless rollouts, which is the modelling choice we attribute to
prior work; noiseaware trains under the true heterogeneous noise and can in
principle discover the reliability weighting from data. The feature set
deliberately includes the interaction between split balance and attribute
identity ($\Hb(p) \times$ attribute one-hot), so a linear model over these
features \emph{can} represent a per-attribute discount on balanced splits,
which is precisely what RAIG encodes. Omitting that interaction would have
guaranteed the baseline could not compete and made any win an artefact of
feature design.

\subsection{Noise Conditions}
\label{sec:conditions}

The simulator draws answer errors from a true per-question crossover vector
$\boldsymbol\varepsilon$; every method sees only its own assumed
$\hat{\boldsymbol\varepsilon}$. The two are separate arguments in the session
signature, and conflating them would make every result circular. Seven
conditions are used throughout.

\begin{enumerate}
\item \emph{Clean}: $\varepsilon \equiv 0$. A sanity condition, not a realistic
one. Its purpose is to verify that every method attains the identifiability
ceiling when the model's noise assumption is the only thing that is wrong.
\item \emph{Homogeneous}: $\varepsilon \equiv 0.098$, the catalogue average of
the tier annotation. Answer noise, but no heterogeneity, hence nothing for a
reliability-aware selector to exploit (Proposition~\ref{prop:uniform}).
\item \emph{Heterogeneous} at $\lambda \in \{0.5, 1.0, 1.5\}$:
$\varepsilon_f = \lambda \cdot \varepsilon_{\mathrm{tier}(f)}$, with
$\lambda=1.0$ the primary condition. $\lambda$ scales user noise without
changing its structure, and locating the value at which reliability awareness
begins to pay is the practitioner-facing question of
Section~\ref{sec:results-lambda}.
\item \emph{Randomised}: the attribute-to-$\varepsilon$ assignment is permuted.
The tier annotation is then uninformative about the truth, and the system is
running on an annotation that has been scrambled behind its back.
\item \emph{Adversarial}: $\varepsilon$ is assigned in inverse rank order of
maximum information gain, so the best-splitting attribute becomes the most
reliable. The annotation is now not merely uninformative but wrong in the
direction most damaging to RAIG.
\end{enumerate}

The last two are the conditions under which the theory predicts our method
should lose, and we report them with the same prominence as the conditions in
which it wins. A method that could not be made to fail by inverting its
assumption would be a method whose assumption did no work.

\subsection{Metrics}
\label{sec:metrics}

Top-1 accuracy at the budget is the headline metric because it is what an
identification system is for. It is also the metric most compressed against the
entropy floor at $T=20$, so we report alongside it: top-5 and top-10 accuracy;
the median rank of the true item; mean reciprocal rank; the negative log
posterior of the true item in bits, which is sensitive to belief quality even
when no prediction is correct; and, for the behavioural claim, the mean true
$\varepsilon$ of the questions each method actually chose to ask, together with
the tier composition of those questions. That last pair is the most direct
evidence for or against the bias, because it measures what the selector did
rather than how well it happened to do.

Ranks use the mid-rank convention, $1 + |\{i : b_i > b^\ast\}| +
(|\{i : b_i = b^\ast\}| - 1)/2$. This matters more than it sounds: early in a
session the belief vector carries very large ties, and the optimistic
convention would credit a method for a tie it has not broken.

\subsection{Statistical Protocol}
\label{sec:stats}

Trials are paired by construction. For each trial index the harness draws one
target item and one stream of uniform variates, and every method under
comparison plays that same target against that same stream
(Algorithm~\ref{alg:eval}). Methods diverge in which questions they ask and
therefore in which crossover probabilities they face, but they consult the same
underlying draws, so a difference in outcome is attributable to question choice
rather than to sampling luck. The primary comparison uses $R=300$ trials at
each of three seeds, $n=900$ paired trials per cell.

Significance is assessed with the exact (binomial) McNemar
test~\cite{mcnemar1947} on the discordant pairs, which is the appropriate test
for paired binary outcomes and does not rely on a normal approximation that
would be poor at these discordance counts. Within each condition the family of
comparisons against the reference method is corrected by the Holm step-down
procedure~\cite{holm1979}. Confidence intervals on accuracies and on paired
differences are bootstrap percentile intervals over $10{,}000$ resamples.

One comparison is designated primary in advance: RAIG-tiered against
IG+soft-uniform, top-1 accuracy, heterogeneous $\lambda=1.0$, $T=20$. Every
other comparison in this paper is exploratory, and we label the tables
accordingly rather than reporting a wall of asterisks and letting the reader
assume they were all planned.
